\begin{abstract}
El banco de baterías es el componente que permite a un sistema fotovoltaico aislado alimentar las cargas durante la noche y en los días sin sol. En este informe se describen las tecnologías de baterías aplicables a un banco residencial de \SI{24}{V}: plomo-ácido de electrolito inundado, plomo-ácido regulada por válvula con separador de vidrio absorbente (AGM) y con electrolito en gel, y ion-litio de fosfato de hierro y litio (LiFePO\textsubscript{4}). Se definen los parámetros de comparación: capacidad nominal, tensión, profundidad de descarga, ciclos de vida, eficiencia, autodescarga, temperatura de operación, mantenimiento y costo. A partir de esta caracterización se identifican opciones comerciales con su ficha técnica y precio, se comparan mediante el costo por kWh útil y se selecciona la batería para el Proyecto~2. \pendiente{completar el resumen con la batería seleccionada y la capacidad del banco cuando se tenga el consumo}
\end{abstract}

\begin{IEEEkeywords}
baterías, sistema fotovoltaico aislado, plomo-ácido, AGM, gel, LiFePO\textsubscript{4}, profundidad de descarga, días de autonomía, análisis comparativo
\end{IEEEkeywords}
